\documentclass{article}

\usepackage[main]{neurips_2026}

\usepackage[utf8]{inputenc}
\usepackage[T1]{fontenc}
\usepackage{hyperref}
\usepackage{url}
\usepackage{booktabs}
\usepackage{amsfonts}
\usepackage{amsmath}
\usepackage{nicefrac}
\usepackage{microtype}
\usepackage{xcolor}
\usepackage{graphicx}
\usepackage{algorithm}
\usepackage{algpseudocode}
\usepackage{multirow}

\title{Function-Aware Fill-in-the-Middle as Mid-Training\\for Coding Agent Foundation Models}

\author{%
  Anonymous Authors \\
  Anonymous Institution \\
  \texttt{anonymous@example.com}
}


\begin{document}
\maketitle

% =============================================================================
% ABSTRACT  (~ 175 words, single paragraph)
% =============================================================================
\begin{abstract}
Coding agents must integrate external tool returns into ongoing reasoning---a
capability that standard left-to-right pretraining on code exposes only in its
forward direction. We observe that the
\textit{action $\rightarrow$ observation $\rightarrow$ continuation} loop of a
coding agent is structurally isomorphic to a function call site, where a
caller binds arguments, a callee returns a value computed elsewhere, and
downstream code consumes that value. This conditioning structure exists at
internet scale in ordinary code. We exploit it through
\emph{function-aware fill-in-the-middle (FIM) mid-training}: a self-supervised
objective that masks functions selected via program dependency graph
analysis and a complexity--inferability double criterion. We mid-train
Qwen2.5-Coder (7B/14B/32B) and Qwen3-8B on a 3B-token decontaminated corpus
drawn from 990 GitHub repositories, then apply existing agentic post-training
pipelines (r2e-gym, swe-smith) without modification. Mid-training improves
SWE-Bench-Verified by $+2.8/+3.0/+3.3$ points at 7B/14B/32B---a consistent
gain across scales---and the improvement holds across post-training pipelines
and base-model families. Beyond in-domain gains, mid-training also mitigates
the capability erosion that agentic post-training otherwise inflicts on
non-agent coding (e.g., LiveCodeBench) and non-coding tool-use benchmarks
(tau-bench, BFCL): although the mid-training corpus contains Python code
only, the function-call inductive bias survives post-training and yields
consistent improvements over post-training alone.
\end{abstract}


% =============================================================================
% 1. INTRODUCTION  (~ 1.3 页)  -- 见 intro.tex
% =============================================================================
\input{sections/intro.tex}
% [Intro content provided separately above; insert here.]


% =============================================================================
% 2. RELATED WORK  (~ 0.6 页)
% =============================================================================
% 四块:
%   2.1 Mid-training / continued pretraining
%       - 引用 MiniCPM, OLMo, Code-Llama long-context mid-training 等
%       - 我们的差异化: agent-oriented mid-training, 用 structured FIM.
%
%   2.2 Fill-in-the-middle pretraining
%       - 引用 Bavarian 2022, Code-Llama, StarCoder, Qwen-Coder, DeepSeek-Coder
%       - 关键: 这些都是 random span FIM, 在 pretraining 阶段大规模混合.
%         我们的差异化: function-aware selection + CoT + mid-training stage.
%   CAUTION: 不要否认前人贡献, frame 成 "extension/specialization".
%
%   2.3 Coding agent foundation models
%       - SWE-Bench, r2e-gym, swe-smith, SWE-agent, OpenHands
%       - Synthetic agent trajectory works (SWE-Gym, Agent-FLAN)
%       - 我们的差异化: 不依赖 trajectory, 用 self-supervised code 信号.
%
%   2.4 Distillation from frontier models [必须有, 主动 disarm]
%       - acknowledge 我们用 Gemini-3 生成 CoT, 引用 Distill-Whisper, Orca,
%         Self-Instruct 等
%       - 强调我们 ablation 拆开 distillation 与 FIM 结构两部分贡献:
%         "We separate the contribution of FIM structure from CoT distillation
%          via a controlled ablation in Section X.X."
% CAUTION: 这一节存在 = 主动 disarm reviewer 的 distillation 质疑.
% =============================================================================
\section{Related Work}
\label{sec:related}

\paragraph{Mid-training and continued pretraining.}
% TODO

\paragraph{Fill-in-the-middle objectives.}
% TODO

\paragraph{Coding agent foundation models.}
% TODO

\paragraph{Distillation from frontier models.}
% TODO (主动 disarm reviewer)


% =============================================================================
% 3. METHOD  (~ 2 页)  -- 已写完, 见 method.tex
% =============================================================================
\input{sections/method.tex}


% =============================================================================
% 4. EXPERIMENTS  (~ 3.0 页, 32B 已完成, 去掉 compute-matched)
% =============================================================================
\section{Experiments}
\label{sec:exp}


% -----------------------------------------------------------------------------
% 4.1 Experimental Setup  (~ 0.4 页)
% -----------------------------------------------------------------------------
% 关键: 必须在 setup 显式说明 evaluation 是在 mid+post (or post-only) 上做的,
% 不是在 mid-only ckpt 上做的. Reviewer 会问. 现在主动说清楚.
%
% Benchmarks 拆三类:
%   * Agent benchmarks (in-domain):    SWE-Bench-Verified, SWE-Bench-Lite
%   * Tool-use benchmarks (OOD):       tau-bench, BFCL
%   * Non-agent coding (regression check):
%       LiveCodeBench, OJBench, FullStackBench-EN/ZH, Terminal-Bench, HumanEval+
%
% Eval protocol:
%   - All evaluations are on the FINAL checkpoint after the full pipeline.
%   - "Baseline" means: base model + post-training (e.g. r2e-gym).
%   - "Ours" means:     base model + FIM mid-training + post-training.
%   - We do NOT evaluate the mid-train-only checkpoint directly because mid-
%     training on FIM data degrades instruction-following ability (Qwen2.5-
%     Coder-Instruct itself acquires IF via standard SFT post-training).
%     Comparing a mid-train-only ckpt to instruction-tuned baselines would
%     be unfair on instruction-following bench. We make this explicit so
%     all "+X" claims are read as "X gained from mid-training that survives
%     subsequent post-training".
%
% Eval scaffold: agent harness, max turns, retrieval setting
% Statistical protocol: 3 seeds, mean ± std, paired bootstrap
%
% Baselines:
%   - Qwen2.5-Coder-Instruct (no agentic training, ceiling for non-agent bench)
%   - r2e-gym  (ours-reproduced + official)
%   - swe-smith (ours-reproduced + official)
%   - +FIM-Midtrain + r2e-gym
%   - +FIM-Midtrain + swe-smith
% -----------------------------------------------------------------------------
\subsection{Setup}
\label{sec:exp:setup}

% TODO benchmarks (3 类) + eval protocol (post-pipeline) + baselines


% -----------------------------------------------------------------------------
% 4.2 Main Results: Agent and Tool-Use Benchmarks  (~ 0.7 页 + 大表)
% -----------------------------------------------------------------------------
% =============================================================================
% TAB 3 [必须, paper 主表 1]: Agent + tool-use 主表
%   行: 跨 7B/14B/32B 用 \multirow 分组, 每组内:
%     - Base instruct (no agentic post-training)
%     - r2e-gym (ours-reproduced)
%     - r2e-gym (official report)         [灰色]
%     - +FIM-Midtrain + r2e-gym           [bold, ours]
%     - swe-smith (ours-reproduced)       [仅 7B 块有]
%     - swe-smith (official report)       [灰色, 仅 7B]
%     - +FIM-Midtrain + swe-smith         [bold, ours, 仅 7B]
%   列: SWE-Bench-V, SWE-Bench-L, tau-bench, BFCL, average
%   每个数字带 ± std (3 seeds).
%
% 数字 (32B 已校正):
%   7B  + r2e-gym:                     SWE-V 15.00, SWE-L 11.33
%   7B  + FIM + r2e-gym:               SWE-V 17.80, SWE-L 15.00
%   7B  + swe-smith:                   SWE-V 12.30, SWE-L 15.20
%   7B  + FIM + swe-smith:             SWE-V 17.60, SWE-L 14.70
%   14B + r2e-gym:                     SWE-V 26.20, SWE-L 18.00
%   14B + FIM + r2e-gym:               SWE-V 29.20, SWE-L 22.00
%   32B + r2e-gym:                     SWE-V 31.80, SWE-L 24.70
%   32B + FIM + r2e-gym:               SWE-V 35.10, SWE-L 26.80
% =============================================================================
%
% 写作 critical paragraph (paper 最强卖点之一):
%   独立段落 highlight cross-domain transfer (tau-bench / BFCL):
%   "Crucially, our mid-training corpus contains exclusively Python code with
%    no tool-use trajectory data. While the absolute tau-bench / BFCL scores
%    of mid-train+post-train models do not exceed the original instruction-
%    tuned base, they consistently exceed those of post-train-only models
%    (+4.0 on tau-bench, +2.5 on BFCL at 14B). Since mid-training contains
%    no tool-use signal, the only mechanism by which it can affect these
%    benchmarks is by installing a structural prior that survives post-
%    training. This directly supports our hypothesis that function calls
%    and agent tool calls share an underlying conditioning structure."
% -----------------------------------------------------------------------------
\subsection{Main Results on Agent and Tool-Use Benchmarks}
\label{sec:exp:main}

% TODO TAB 3 (含 32B 完整数据)
% TODO 主结果段落
% TODO 单独一段 cross-domain transfer (tau-bench / BFCL)


% -----------------------------------------------------------------------------
% 4.3 Capability Preservation on Non-Agent Coding  (~ 0.6 页 + 表)
%     [REWRITTEN: 强化"防退化"叙事]
% -----------------------------------------------------------------------------
% 这一节承担了原 compute-matched 的 selling point.
% 核心 message: 单独 r2e-gym post-training 在 non-agent coding bench 上严重
% 退化, 加 mid-training 后退化被基本消除, 部分 benchmark 还有提升.
% --> mid-training 改善了 agent post-training 的 capability cost-benefit profile.
%
% =============================================================================
% TAB 4 [必须, paper 主表 2]: Non-agent coding 主表 (14B 完整对照)
%   行结构:
%     Qwen2.5-Coder-14B-Instruct          (no agentic post-training, ceiling)
%     14B + r2e-gym                        (post-training only, 严重退化)
%     14B + FIM-Midtrain + r2e-gym         (我们的, 接近恢复)
%   列: LiveCodeBench, OJBench, FullStackBench-EN, FullStackBench-ZH,
%       Terminal-Bench, HumanEval+, average
%
% 数字 (14B):
%   Instruct:               LiveCode 37.2, OJBench 5.20, FSB-EN 53.8,
%                           Terminal 0,    tau 5.70,     BFCL 23.2  (avg 20.85)
%   + r2e-gym:              LiveCode 24.1, OJBench 2.80, FSB-EN 47.72,
%                           Terminal 1.28, tau 0.40,     BFCL 15.7  (avg 15.33)
%   + FIM + r2e-gym:        LiveCode 35.2, OJBench 4.74, FSB-EN 48.07,
%                           Terminal 2.47, tau 4.40,     BFCL 18.2  (avg 18.85)
% =============================================================================
%
% 写作要点 (3 段):
%   (a) Capability regression as the hidden cost of agentic post-training.
%       单独 r2e-gym 让 14B LiveCode 从 37.2 砸到 24.1, FSB-EN 从 53.8 到 47.72.
%       这是现有 agent 工作很少 report 的代价.
%   (b) Mid-training 大幅消除该退化:
%       LiveCode 35.2 (恢复 11.1 of 13.1 drop),
%       OJBench 4.74 (vs ceiling 5.20),
%       FSB-EN 48.07 (recovery 0.35).
%       Average 18.85 vs post-only 15.33: +3.5 跨 6 个 benchmark.
%   (c) 跨域泛化 (训练数据是 Python, 多/异构 任务也受益):
%       - Terminal-Bench 1.28 -> 2.47 (closest to non-coding tool use)
%       - FullStackBench-ZH 提升 (语言泛化)
%       - 这些点支持 "structural prior 而非 task-specific shortcut".
%
% CAUTION:
%   - Frame 成 "mid-training fixes a side-effect of post-training while still
%     improving in-domain", 不是 "we beat post-training".
%   - 由于 compute, 此 controlled regression 分析仅在 14B; in-domain SWE 一致
%     的 7B/32B 趋势 suggests 类似 pattern, 但留作 limitation.
% -----------------------------------------------------------------------------
\subsection{Capability Preservation on Non-Agent Coding}
\label{sec:exp:gen}

% TODO TAB 4 (14B 三行对照)
% TODO 段落 (a): post-training 单独导致 capability regression
% TODO 段落 (b): mid-training 缓解退化, 6-bench average +3.5
% TODO 段落 (c): 跨域 (Terminal-Bench, FSB-ZH) 进一步支持 structural prior 论点


% -----------------------------------------------------------------------------
% 4.4 Cross-Base-Model Validation: Qwen3-8B  (~ 0.3 页 + 小表)
% -----------------------------------------------------------------------------
% 支持 abstract/intro 的 cross-base-model claim.
%
% =============================================================================
% TAB 5 [必须]: Qwen3-8B 小表
%   行: Qwen3-8B + swe-lego (50%)                  baseline
%       Qwen3-8B + FIM-Midtrain + swe-lego (50%)   [bold, 数字标红]
%   列: SWE-Bench-V (其他列等结果再补)
%
% 数字 (待更新, 标红):
%   baseline:           SWE-V 30.2
%   + FIM-Midtrain:     SWE-V 33.4   (+3.2)
% =============================================================================
%
% 写作要点:
%   - "to verify the approach is not specific to Qwen2.5-Coder, we repeat the
%     core experiment on Qwen3-8B"
%   - Honest about pipeline: 这里使用 swe-lego (50% data) 做 post-training,
%     与 Qwen2.5-Coder 的 r2e-gym 不同, 因此 cross-base 这一变量同时伴随
%     pipeline 变化; 我们 frame 成 "在另一组合下仍观察到类似量级的提升 (+3.2)"
%   - 数字标红, paper 投稿前更新 + 去掉颜色
%
% CAUTION:
%   - 措辞谨慎: 一组 base model 不能 generalise 到 "all base model"
%   - Limitation 重申一次
% -----------------------------------------------------------------------------
\subsection{Cross-Base-Model Validation}
\label{sec:exp:crossbase}

% TODO TAB 5: Qwen3-8B (数字标红, 待更新)
% TODO 段落, 措辞谨慎


% -----------------------------------------------------------------------------
% 4.5 Consistency Across Model Sizes  (~ 0.3 页 + 1 张图)
%     [UPDATED: 32B 校正后, gain consistent rather than growing]
% -----------------------------------------------------------------------------
% =============================================================================
% FIG 6 [必须]: 三 size scaling 图.
%   X 轴 model size (7B/14B/32B, log scale),
%   Y 轴 SWE-Bench-Verified,
%   两条线: r2e-gym only, +FIM-Midtrain + r2e-gym.
%   每个 size 旁标注 absolute gain (+2.8 / +3.0 / +3.3).
%
% 数字 (校正后):
%   7B:  15.00 -> 17.80   gain +2.8
%   14B: 26.20 -> 29.20   gain +3.0
%   32B: 31.80 -> 35.10   gain +3.3
% =============================================================================
%
% 写作 framing (低调, consistent 而非 growing):
%   "Across the three sizes evaluated, the absolute gain from mid-training
%    remains within a narrow band (+2.8 to +3.3) and shows no sign of
%    diminishing with scale. This addresses a natural concern that larger
%    models, having seen more pretraining tokens, would already have
%    absorbed the structural prior; in our experiments the gain at 32B is
%    on par with that at 7B and 14B. The relative gain is largest at 7B
%    (+18.7%) and smaller at 14B/32B (+11.5%/+10.4%), as expected from a
%    higher baseline. Whether the gain remains consistent at 70B+ is left
%    to future work."
%
% CAUTION:
%   - 不 claim "scales", 用 "consistent across the scales evaluated"
%   - 主动 acknowledge relative gain 在 14B 略低于 7B, 用 "higher baseline"
%     解释, 避免 reviewer 拿这点攻击
%   - 不外推到 70B+
% -----------------------------------------------------------------------------
\subsection{Consistency Across Model Sizes}
\label{sec:exp:scaling}

% TODO FIG 6: scaling curve (三个 size 实测点)
% TODO 段落: consistent gain across scales, 反驳 "大模型已经会了"


% -----------------------------------------------------------------------------
% [REMOVED] Compute-Matched Comparison
%
% RATIONALE: mid-training 与 post-training 不必在具体下游 SWE 任务上直接做
% compute-match. Section 4.3 已经显示 post-training 单独对 non-agent coding
% 与 tool-use 能力造成显著退化 (LiveCode -13.1, BFCL -7.5, tau-bench -5.3),
% 而 mid-training 能在保留这些能力的同时提升 in-domain 性能, 这本身就是比
% compute-matched curve 更直接的 "投资 mid-training 而非更多 post-training
% data" 的论据.
% -----------------------------------------------------------------------------


% -----------------------------------------------------------------------------
% 4.6 Ablation Studies  (~ 0.7 页, 表 + 文字 + 图)
% -----------------------------------------------------------------------------
% Ablation 维度 (全部在 14B + r2e-gym 上 controlled):
%
%   (A) CoT source [critical, disarm distillation 质疑]:
%       - Baseline (no mid-train, post-train only)
%       - FIM mid-train, no CoT
%       - FIM mid-train, 14B-self CoT
%       - FIM mid-train, Gemini-3 CoT  (full method)
%       Narrative:
%         "Removing CoT reduces SWE-V from X to Y but still beats no
%          mid-training by Z, indicating FIM structure contributes
%          independently of distillation signal. Replacing Gemini-3 with
%          self-generated CoT recovers most but not all of the gain,
%          suggesting that CoT quality matters but is not the sole driver."
%       CAUTION: 主动写 "due to compute constraints, we conduct this
%                ablation on the 14B configuration".
%
%   (B) Selection algorithm:
%       - Random function selection
%       - PDG only (no inferability score)
%       - Inferability only (no PDG)
%       - Full (PDG + complexity + inferability)
%
%   (C) Mask granularity (1 / 2 / 3 connected functions / mixed)
%
%   (D) Score threshold sensitivity (tau_d, tau_FIM sweep)
%
%   (E) Mid-training token budget scaling (0.5B / 1B / 2B / 3B / 5B)
%       回答 "more is better?" + saturation 行为
%
% =============================================================================
% TAB 6 [必须]: 综合 ablation 表
%   一张大表, 五块 sub-table 或 multirow:
%   (A) CoT source: 4 行
%   (B) Selection: 4 行
%   (C) Granularity: 4 行
%   (D) Threshold: 5 行 sweep
%   (E) Token budget: 5 行 sweep
%   列: SWE-Bench-V (主指标), 可选 SWE-Bench-L
%
% FIG 9 [建议]: token budget scaling curve
%   X 轴 mid-training tokens (0.5B / 1B / 2B / 3B / 5B), log scale,
%   Y 轴 SWE-Bench-Verified.
% =============================================================================
% -----------------------------------------------------------------------------
\subsection{Ablation Studies}
\label{sec:exp:ablation}

% TODO TAB 6: 综合 ablation 表
% TODO CoT source ablation (4 配置)
% TODO selection algorithm ablation
% TODO mask granularity ablation
% TODO threshold sensitivity sweep
% TODO 实验 (强烈建议): mid-training token budget scaling
% TODO FIG 9: token budget scaling curve


% =============================================================================
% 5. ANALYSIS  (~ 0.9 页)
% =============================================================================
\section{Analysis}
\label{sec:analysis}


% -----------------------------------------------------------------------------
% 5.1 Behavioral Analysis  (~ 0.4 页)
% -----------------------------------------------------------------------------
% Trajectory-level metrics, mid-training 前后对比:
%   - Average tool calls per task
%   - Recovery rate after error / unexpected output
%   - Average context usage at success
%   - max-turn 限制下的 success curve
%
% =============================================================================
% TAB 7 [建议]: trajectory metrics 对比 (14B)
%   行: r2e-gym only / +FIM + r2e-gym
%   列: avg tool calls, recovery rate, context @ success, success @ max-turn=k
% =============================================================================
% -----------------------------------------------------------------------------
\subsection{Behavioral Changes in Agent Trajectories}
\label{sec:analysis:behavior}

% TODO TAB 7: trajectory-level metrics
% TODO 实验: parse SWE-agent log 提取 metrics


% -----------------------------------------------------------------------------
% 5.2 Failure Mode Analysis  (~ 0.3 页)
% -----------------------------------------------------------------------------
% =============================================================================
% FIG 8 [建议]: stacked bar chart, 失败类型分布对比
%   X 轴: r2e-gym only vs +FIM + r2e-gym
%   Y 轴: # failed tasks, 按 failure type stacked
%   Failure types: localization error, patch error, env error, infinite loop,
%                  timeout, etc.
% =============================================================================
%
% 写作 takeaway: 提升集中在 "需要 function-level reasoning across files"
% 的任务, 对应 motivation.
% -----------------------------------------------------------------------------
\subsection{Where Does Mid-Training Help Most?}
\label{sec:analysis:failure}

% TODO FIG 8: failure mode breakdown
% TODO 分析: 标注 SWE-Bench task 类型


% -----------------------------------------------------------------------------
% 5.3 General Capability Forgetting Check  (~ 0.2 页)
% -----------------------------------------------------------------------------
% 这一节与 Section 4.3 互补: 4.3 看 coding/tool-use 任务退化,
% 这里看更广义的通用能力 (语言/数学/IF).
%
% =============================================================================
% TAB 8 [必须]: 通用 benchmark forgetting check (14B)
%   行:
%     Qwen2.5-Coder-14B-Instruct       baseline ceiling
%     14B + r2e-gym                     post-training only
%     14B + FIM + r2e-gym               ours
%   列: MMLU, MATH, IFEval, average
% =============================================================================
%
% CAUTION: 不测的话 reviewer 必然假设有 forgetting 问题.
%          主动 report 比被发现强.
% -----------------------------------------------------------------------------
\subsection{Preservation of General Capabilities}
\label{sec:analysis:forgetting}

% TODO TAB 8: forgetting check
% TODO 实验: MMLU / MATH / IFEval 14B baseline vs +FIM


% =============================================================================
% 6. LIMITATIONS AND DISCUSSION  (~ 0.4 页)
% =============================================================================
% 必须 acknowledge:
%   1. Python-heavy benchmark; 跨语言 generalisation 仅由 FullStackBench-ZH
%      间接支持, 未在主表覆盖 Java/C++/Rust 等.
%   2. CoT 监督依赖 frontier model API (Gemini-3-Preview); 自蒸馏 ablation
%      恢复了大部分 gain (Section 4.6), 但完全开源复现需要相应替代.
%   3. PDG 分析对动态语言不完美 (动态 dispatch, decorator, monkey patching);
%      详见 Appendix.
%   4. Cross-base-model 仅在 Qwen3-8B + swe-lego 上做, 该实验同时伴随
%      post-training pipeline 变化, 因此 cross-family 结论的强度受限.
%      未覆盖 Llama, DeepSeek 等其他 family.
%   5. Capability preservation 的 controlled comparison 仅在 14B (Section 4.3);
%      7B/32B 的 in-domain SWE-Bench 趋势 suggests 类似 pattern 但未 controlled.
%   6. Function-aware FIM 假设代码 modular; 极端 monolithic 或 inline-heavy
%      代码上的 selection pipeline 行为未充分研究.
% =============================================================================
\section{Limitations and Discussion}
\label{sec:limitations}

% TODO: 6 条 limitations


% =============================================================================
% 7. CONCLUSION  (~ 0.25 页)
% =============================================================================
% 4-5 句:
%   - 重申 framing (function call ↔ agent step 的 structural isomorphism)
%   - 重申 method (function-aware FIM mid-training, PDG + 双信号 + CoT)
%   - 重申结果: 三轴 robustness 一致 + 同一 corpus 跨任务族提升
%   - future direction: 跨语言, 跨 base family 扩展, mid-training 与
%     RL post-training 的协同
% =============================================================================
\section{Conclusion}
\label{sec:conclusion}

% TODO: 4-5 句 conclusion


% =============================================================================
% References & Appendix
% =============================================================================
\bibliographystyle{plainnat}
% \bibliography{refs}

\section*{References}
% TODO: 整理参考文献
% 必引文献 cluster:
%   FIM:           Bavarian 2022, Code-Llama, StarCoder, Qwen-Coder,
%                  DeepSeek-Coder
%   Mid-training:  MiniCPM, OLMo, 任何 explicit "mid-training" paper
%   Coding agent:  SWE-Bench, SWE-agent, OpenHands, Aider, r2e-gym, swe-smith,
%                  swe-lego, R2E, SWE-Gym
%   Distillation:  Orca, Self-Instruct, Distill-Whisper
%   Tool use eval: tau-bench, BFCL
%   General eval:  MMLU, MATH, IFEval (forgetting check)
%   Program analysis: PDG classical refs

\appendix

% =============================================================================
% APPENDIX
%
%   A. Detailed dataset statistics (含 78K filtered files, 33M lines, 平均 LoC,
%      10 类 category 分布, license 分布, decontamination 列表)
%   B. Hyperparameters (含 mid-training 3B token budget 详细 schedule, LR sweep)
%   C. Algorithm details
%      - PDG 实现细节 (Python AST 处理动态特性的 workaround)
%      - 复杂度/可推断性 score 完整公式与超参
%      - Group topology 完整 8 种 pattern 定义
%      - CoT prompt template (Gemini-3 用什么 prompt)
%      - 14B-self CoT prompt template (与 Gemini-3 一致, 仅换模型)
%   D. Additional results
%      - 每个 benchmark 的 per-task breakdown
%      - Qwen3-8B 上更详细的 trajectory 分析 (如有)
%      - 7B / 32B 上的 capability preservation sanity check (如时间允许)
%   E. Qualitative examples (一两个 SWE-Bench task, 对比 baseline 与 ours
%      的 trajectory)
%   F. Reproducibility checklist
% =============================================================================
\section{Dataset Details}
\label{app:data}

\section{Training Hyperparameters}
\label{app:hp}

\section{Algorithm Details}
\label{app:algo}

\section{Additional Results}
\label{app:results}

\section{Qualitative Examples}
\label{app:examples}

\input{checklist.tex}

\end{document}